% Paste-ready RAGAS additions for the dissertation.

\subsection{RAGAS-Assisted Evaluation}

In addition to the manual scoring rubric, the final evaluation incorporates RAGAS as a secondary analysis layer for the comparative benchmark. The motivation for adding RAGAS is not to replace the project-specific rule-fidelity checks, but to provide a standardised retrieval-evaluation framework that can be compared with the manual findings. RAGAS is explicitly designed to evaluate retrieval-augmented generation pipelines using metrics that separate response quality from retrieval quality, which is useful for determining whether failures arise from poor retrieval, poor use of retrieved evidence, or broader weaknesses in structured generation.

The benchmark outputs are converted into a RAGAS-compatible dataset containing four core fields: the user input, the generated response, the retrieved contexts, and a reference answer. For this project, the user input is the structured generation request, the response is a normalised textual rendering of the generated JSON output, the retrieved contexts are the rule passages returned by the retriever, and the reference answer is a curated benchmark reference description for each prompt. This allows the same 10-prompt benchmark used for manual scoring to be analysed with additional automatic metrics.

The RAGAS metrics selected for this project are:
\begin{itemize}
    \item \textbf{Faithfulness:} whether the generated response is supported by the retrieved rule context.
    \item \textbf{Answer relevancy:} whether the generated response addresses the user request.
    \item \textbf{Context precision:} whether the retrieved context is relevant to producing the answer.
    \item \textbf{Context recall:} whether the retrieved context contains the information needed to support the answer.
    \item \textbf{Answer correctness:} whether the generated response matches the benchmark reference description.
\end{itemize}

These metrics are used only as a complement to the dissertation's manual rubric. This limitation is important. RAGAS can indicate whether retrieved evidence appears relevant and whether a response is broadly aligned with a reference answer, but it does not directly guarantee D\&D 5e legality. A response may still score reasonably on relevancy while violating spell progression, proficiency constraints, or derived-stat rules. For that reason, the project retains manual \textit{Rules fit} scoring and deterministic validation as the primary measures of mechanical correctness.

\subsection{How RAGAS Is Used in the Comparative Study}

The same three benchmark conditions are used for both manual scoring and RAGAS scoring:
\begin{itemize}
    \item \textbf{No-RAG}
    \item \textbf{RAG}
    \item \textbf{Validated iterative RAG}
\end{itemize}

For each of the 30 outputs in the benchmark, the system records the generated answer, the parsed structured output, the retrieved sources, the validation issues, and the model used. The RAGAS evaluation script then consumes these saved artifacts and produces condition-level summaries. This enables direct comparison between the manual results and the RAGAS metrics.

The dissertation uses this comparison in two ways. First, it checks whether improvements in manual \textit{Rules fit} are mirrored by improvements in faithfulness or context precision. Second, it highlights cases where they diverge. Such divergence is analytically useful because it can indicate that retrieval is working while downstream structured reasoning is still failing, or conversely that retrieval quality is weak even when the final output appears superficially plausible.

\subsection{Interpretive Limits of RAGAS in This Project}

RAGAS was originally developed for more conventional retrieval-augmented generation tasks, particularly tasks closer to grounded question answering. D\&D character generation differs from those settings because the output is not a short textual answer but a structured build with multiple interacting constraints. This makes RAGAS informative but insufficient on its own.

Accordingly, the dissertation treats RAGAS as a diagnostic layer rather than the sole evaluation method. If manual rule-fidelity scores improve while RAGAS faithfulness remains flat, this suggests that deterministic validation is correcting issues that the retrieval metrics do not capture directly. If RAGAS context precision improves without a corresponding improvement in \textit{Rules fit}, this suggests that retrieval alone is not enough and that the model is still misapplying valid context. These are precisely the distinctions that the comparative ablation is intended to surface.
